\documentclass[11pt]{article}
\usepackage{fontspec}
\usepackage[french,english]{babel}
\usepackage{amsmath}
\usepackage{amssymb}
\usepackage{amsthm}
\usepackage{mathrsfs}
\usepackage{tikz-cd}
\usetikzlibrary{babel}
\usepackage{enumitem}

\newcommand{\printedpage}[1]{\marginpar{\footnotesize p.~#1}}
\newcommand{\msp}[1]{\marginpar{\footnotesize [#1]}}
\newcommand{\Fl}{\mathrm{Fl}}
\newcommand{\M}{\mathcal{M}}

\begin{document}

\printedpage{14}
% head: A. Grothendieck : Dérivateurs, Ch. XVII, Édition des Universités de Montpellier II et Paris VII

\msp{page 24}

\textbf{Lemma.} \textit{Let $P \subset I$ be the set of $i_0 \in I$ such that for every open subset $K$ of $J = I/i_0$, the arrow $u_{K,J}$ is in $C$. This said, let $i_0 \in I$ be such that every $i < i_0$ is in $P$. Then $i_0 \in P$.} (We assume $u_\bullet \in C_I$, but $I$ arbitrary.)

Indeed, let $K \subset J = I/i_0$ be an open subset. If $K = J$, it is trivial that $u_{K,J} \in C$ (it is an isomorphism). If $K = J$, we have % [sic] Si K = J (second occurrence) -> Si K \neq J, the first case already treats K = J
\[
K \ \subset\ J' \overset{\text{def}}{=} J\backslash\{i_0\} \ \subset\ J,
\]
and to prove that $u_{K,J} \in C$, it suffices to prove that $u_{K,J'}, u_{J',J} \in C$. Now $u_{J',J} \in C$ by the hypothesis $u_\bullet \in C_I$. And $u_{K,J'} \in C$ by what has been proved for the case of $J'$, taking into account that by the hypothesis on $i_0$, every $i \in J'$ (i.e. every $i < i_0$) is in $P$. % [sic] par ce que a été prouvé -> par ce qui a été prouvé, grammar

\textbf{Corollary.} \textit{If $I$ is artinian, $P = I$.}

For any subset $P$ of $I$ having the property just stated ($I/i_0 \subset P \Longrightarrow i_0 \in P$) is equal to $I$. If indeed there were $x_0 \notin P$, there would exist $x_1 < x_0$ with $x_0 \notin P$, and by induction we would construct a strictly decreasing sequence in $I$, absurd. % [sic] avec x_0 \notin P -> avec x_1 \notin P, the new element must lie outside P

\msp{page 25}

\section*{2\quad Description of $C_I \subset \Fl(\M(I^\circ))$, for $I$ an ordered set}

Let $\M$ be a category, stable under (small) $\varinjlim$ , $C \subset \Fl(\M)$ (``cofibrations''). We assume
\[
\left\{
\begin{array}{lll}
\text{a)} & C & \text{contains the isomorphisms and is stable}\\
 & & \text{under cobase change.}\\[1ex]
\text{b)} & C & \text{stable under transfinite composition.}
\end{array}
\right.
\tag*{(2.1)}
\]

Let $I$ be a (small) ordered set, and consider
\[
\M(I^\circ) \ =\ \underline{\mathrm{Hom}}(I,\M),
\]
we are going to define a subset
\[
C_I \ \subset\ \Fl(\M(I^\circ))
\]
as follows, for $u_\bullet : X_\bullet \longrightarrow Y_\bullet$ an arrow in $\M(I^\circ)$ :
\[
u_\bullet \in C_I \quad\Longleftrightarrow\quad
\begin{array}{l}
\forall\ i_0 \in I \text{ and for every open subset } J \subset I/i_0 = K,\\
\text{the arrow } u_{J,K} \text{ is in } C.
\end{array}
\tag*{(2.2)}
\]

Here is the definition of $u_{J,K}$. For every open subset $J$ of $I$, we set
\[
\left\{
\begin{array}{rclrcl}
X_J & = & \varinjlim\limits_{i\in J} X_i, & Y_J & = & \varinjlim\limits_{i\in J} Y_i,\\[2ex]
u_J : X_J & \longrightarrow & Y_J, & u_J & = & \varinjlim\limits_{i\in J} u_i,
\end{array}
\right.
\tag*{(2.3)}
\]
and we consider, for $K \subset J$, the diagram
% foot: Version 9 février 2019

\printedpage{15}
% head: A. Grothendieck : Dérivateurs, Ch. XVII, Édition des Universités de Montpellier II et Paris VII

\msp{page 26}

\[
\begin{tikzcd}
X_K \arrow[r, "\tau^X_{K,J}"] \arrow[d, "u_K"'] \arrow[dr, phantom, "\text{\scriptsize cocart.}"] & X_J \arrow[d, "\tilde{u}_{K,J}"'] \arrow[dd, bend left=40, "u_J"] \\
Y_K \arrow[r, "\tau^Y_{K,J}"] & Y_{K,J} \rlap{$\ = Y_K \amalg_{X_K} X_J$} \arrow[d, "u_{K,J}"'] \\
 & Y_J \ .
\end{tikzcd}
\]

\textbf{Theorem 2.1.} \textit{Let $\M$ be a $\mathfrak{U}$-category stable under small $\varinjlim$ , $C \subset \Fl(\M)$, satisfying the stability conditions} \textup{(2.1)} \textit{(under cobase change and under transfinite composition), $I$ a small ordered set, $u_\bullet : X_\bullet \longrightarrow Y_\bullet$ an arrow in $\M(I^\circ)$.}
\begin{itemize}[leftmargin=3em]
\item[$1^\circ$)] \textit{We assume $u_\bullet \in C_I$ (cf.} \textup{(2.2)}\textit{). Let $K \subset J$ be two open subsets of $I$. Then $u_{K,J}$ (defined in} \textup{(2.4)}\textit{) is in $C$.}
\item[$2^\circ$)] \textit{Suppose $I$ is artinian. In order to have $u_\bullet \in C_I$, i.e. for $u_{K,J}$ to be in $C$ for every pair $(K \subset J)$ of open subsets such that $J$ has a greatest element, i.e. is of the form $I/i_0$ (i.e. also, by} $1^\circ$\textit{), for} \textup{every} \textit{pair of open subsets $J \subset K$), (it is necessary and) it suffices that this be so in the particular case where $J = I/i_0$, $J' = I'/i_0 \overset{\text{def}}{=} I/i_0 - \{i_0\} = \{i \in I \mid i < i_0\}$.}

\msp{page 27}
\item[$3^\circ$)] \textit{Suppose again $I$ is artinian}\footnote[11]{I do not know whether this hypothesis is necessary, cf. comment p. 39.}\textit{. Let $W \subset \Fl(\M)$ be a mildly saturated localizer, such that $TC \overset{\text{def}}{=} C \cap W$ also satisfies the stability conditions} \textup{(2.1)} \textit{(under cobase change, under transfinite composition). The following conditions on an arrow $u_\bullet : X_\bullet \longrightarrow Y_\bullet$ in $\M(I^\circ)$ are equivalent.}
\begin{itemize}[leftmargin=2em]
\item[a)] \textit{$u_\bullet \in (TC)_I$, i.e. for every $i_0 \in I$ and $K$ an open subset of $J = I/i_0$, the arrow $u_{K,J}$ is in $TC$.}
\item[b)] \textit{$u_\bullet \in C_I \cap W_{I^\circ}$, i.e. for every $i_0 \in I$, we have $u_{i_0} \in W$, and for every open subset $K$ of $I/i_0 = J$, we have $u_{K,J} \in C$.}
\end{itemize}
\end{itemize}

Let us give the main corollaries right away.

\textbf{Corollary 2.2.} \textit{Suppose $(u_\bullet : X_\bullet \longrightarrow Y_\bullet) \in C_I$. Then}
\begin{itemize}[leftmargin=3em]
\item[a)] \textit{$\forall\ i \in I$, $u_i : X_i \longrightarrow Y_i$ is in $C$.}
\item[b)] \textit{The arrow}
\[
u_I \ :\ X_I \longrightarrow Y_I
\tag*{(2.2.1)}
\]
\textit{is in $C$. More generally, for every open subset $J$ of $I$, the arrow}
\[
u_J \ :\ X_J \longrightarrow Y_J
\tag*{(2.2.2)}
\]
\textit{is in $C$.}
\end{itemize}
% foot: Version 9 février 2019

\printedpage{16}
% head: A. Grothendieck : Dérivateurs, Ch. XVII, Édition des Universités de Montpellier II et Paris VII

\msp{page 28}

For a), we apply $1^\circ$) (or simply the \textit{definition} of $C_I$) in the case of the inclusion of open subsets $\emptyset \hookrightarrow I/i$, since we have
\[
u_{\emptyset,I/i} \ \simeq\ u_i \ .
\]
For b), we apply $1^\circ$) to the inclusion $\emptyset \longrightarrow I$, since
\[
u_{\emptyset,I} \ \simeq\ u_I \ ,
\]
and more generally
\[
u_{\emptyset,J} \ \simeq\ u_J \ .
\]

\textbf{Corollary 2.3.} \textit{Suppose $u_\bullet \in C_I \cap W_{I^\circ}$. Then for every inclusion of open subsets $K \subset J$ of $I$, we have $u_{K,J} \in TC$, and in particular (taking $K = \emptyset$), we have $u_J \in TC$. Finally, for every $i \in I$, we have $u_i \in TC$.}

\textbf{Corollary 2.4.} \textit{Let $X_\bullet$ in $\M(I^\circ)_{\mathrm{c}}$ be a $C$-cofibrant direct system, i.e. such that $\emptyset_I \longrightarrow X_\bullet$ is in $C_I$, i.e. such that for every $i_0 \in I$ and every open subset $K$ of $J = I/i_0$, the transition arrow $\tau_{K,J} : X_K = \varinjlim\limits_{i\in K} X_i \longrightarrow X_J \simeq X_{i_0}$ is in $C$. Then for} \textup{every} \textit{pair of open subsets $K \subset J$, the arrow}
\[
\tau_{K,J} \ :\ \varinjlim_{i\in K} X_i \longrightarrow \varinjlim_{i\in J} X_i
\]
\textit{is in $C$}\footnote[12]{and conversely, characterizes the $C$-cofibrant objects.}\textit{. In particular (taking $K = I/k$, $J = I/j$ for $k \leq j$), if $k \leq j$ are in $I$, the transition morphism}
\[
\tau_{k,j} \ :\ X_k \longrightarrow X_j
\]
\msp{page 29}
\textit{is in $C$. Finally, if $W \subset \Fl(\M)$ is as in} \textup{(2.1.\ $3^\circ$)}\textit{, and if $\emptyset_I \longrightarrow X$ is in $W_{I^\circ}$, i.e. the $X_i$ are $W$-coaspheric, then $X_\bullet$ is in $\M(I^\circ)_{TC}$ (where $TC = C \cap W$), and for every open subset $J$ of $I$, $X_J = \varinjlim\limits_{i\in J} X_i$ is in $\M_{TC}$, hence $X_J$ is $W$-coaspheric, and for every pair of open subsets $K \subset J$, the transition arrow $\tau_{K,J}$ is in $TC$.}

Let us point out finally the following complement to theorem 2.1, which will follow from its proof :

\textbf{Corollary 2.5.} \textit{Suppose that the topological space associated with $I$ (whose open subsets are the open subsets of $I$, and the $I/i$ form a basis of open subsets) is} \textup{noetherian}\textit{, i.e. every increasing sequence of open subsets is stationary. Then all the conclusions in theorem} \textup{2.1} \textit{and its corollaries remain valid, assuming only the stability of $C$ or of $TC = C \cap W$ under simple composition (instead of transfinite composition), in addition to the other two stabilities.}
% foot: Version 9 février 2019

\printedpage{17}
% head: A. Grothendieck : Dérivateurs, Ch. XVII, Édition des Universités de Montpellier II et Paris VII

\msp{page 30}

\textsc{Proof} of theorem 2.1 and of corollary 2.5.

2.1. $1^\circ$) We may assume $J = I$, $K \subset I$. We construct a strictly increasing transfinite sequence of open subsets $I_\alpha$, with
\[
\left\{
\setlength{\arraycolsep}{4pt}
\begin{array}{rcll}
I_0 & = & K, & \\[1ex]
I_{\alpha+1} & = & I_\alpha \cup I/s_\alpha & \text{where } s_\alpha \in I\backslash I_\alpha \text{ (if } I_\alpha \neq I\text{, otherwise, the transfinite construction stops),}\\[1ex]
I_\alpha & = & \bigcup_{\beta<\alpha} I_\beta & \text{if } \alpha \text{ a limit ordinal.}
\end{array}
\right.
\]

There exists an $\alpha_0$ such that we have $I_{\alpha_0} = I$. We then see as in paragraph 1
\msp{page 31}
that for every $\alpha$, we have
\[
u_{I_\alpha,I_{\alpha+1}} \ \in\ C,
\tag*{$(*)$}
\]
for it is deduced from $u_{(I/s_\alpha)\cap I_\alpha,I/s_\alpha}$ by cobase change. Cf. diagram (6), p. 17, which ought to be extracted as a lemma, stripped of the technical peculiarities of the situation, i.e. a lemma on homomorphisms of cartesian squares :

\textbf{Lemma 2.6.} \textit{Let there be a homomorphism of cocartesian squares $X_\bullet$, $Y_\bullet$, in a category $\M$ with amalgamated sums (\textbf{N.B.} it suffices that the amalgamated sums occurring above exist) :}
\[
\begin{tikzcd}[row sep=small, column sep=small]
X_0 \arrow[rr] \arrow[dr] \arrow[dd, "f_0"'] & & X_1 \arrow[dr] \arrow[dd, "f_1" near end] & \\
 & X'_0 \arrow[rr, crossing over] & & X'_1 \arrow[dd, "f'_1"] \\
Y_0 \arrow[rr] \arrow[dr] & & Y_1 \arrow[dr] & \\
 & Y'_0 \arrow[rr] \arrow[from=uu, crossing over, "f'_0" near start] & & Y'_1 \ .
\end{tikzcd}
\tag*{(2.6.1)}
\]

\textit{Let $Z_0$, $Z_1$ be the two lateral squares and consider, for every commutative square}
$Z = \begin{tikzcd}[ampersand replacement=\&, baseline=(\tikzcdmatrixname.center)]
X \arrow[r] \arrow[d] \& X' \arrow[d] \\
Y \arrow[r] \& Y'
\end{tikzcd}$ \textit{, the arrow $u_Z$ of the diagram}
\[
\begin{tikzcd}
X \arrow[r] \arrow[d] \arrow[dr, phantom, "\text{\scriptsize cocart.}"] & X' \arrow[d] & \\
Y \arrow[r] & \tilde{Z} \arrow[r, "u_Z"'] & Y' \ .
\end{tikzcd}
\]
% foot: Version 9 février 2019

\printedpage{18}
% head: A. Grothendieck : Dérivateurs, Ch. XVII, Édition des Universités de Montpellier II et Paris VII

\msp{page 32}

\textit{Then the arrow $u_{Z_1}$ is deduced from $u_{Z_0}$ by cobase change, more precisely, we have a cocartesian square, where} $T = \begin{tikzcd}[ampersand replacement=\&, baseline=(\tikzcdmatrixname.center)]
X_0 \arrow[r] \arrow[d] \& X_1 \arrow[d] \\
Y_0 \arrow[r] \& Y_1
\end{tikzcd}$ \textit{and} $T' = \begin{tikzcd}[ampersand replacement=\&, baseline=(\tikzcdmatrixname.center)]
X'_0 \arrow[r] \arrow[d] \& X'_1 \arrow[d] \\
Y'_0 \arrow[r] \& Y'_1
\end{tikzcd}$ \textit{denote the back and front faces of the cube} \textup{(2.6.1) :}
\[
\begin{tikzcd}
\tilde{Z}_0 \arrow[r, "u_{Z_0}"] \arrow[d] \arrow[dr, phantom, "\text{\scriptsize cocart.}"] & Y'_0 \arrow[d] \\
\tilde{Z}_1 \arrow[r, "u_{Z_1}"] & Y'_1 \ ,
\end{tikzcd}
\tag*{$(*)$}
\]
\textit{which we see on the base of the diagram}
\[
\begin{tikzcd}[row sep=small, column sep=small]
X_0 \arrow[rr] \arrow[dr] \arrow[dd, "f_0"'] & & X_1 \arrow[dr] \arrow[dd, "f_1" near end] & & \\
 & X'_0 \arrow[rr, crossing over] & & X'_1 \arrow[dd] \arrow[dddr, "f'_1"] & \\
Y_0 \arrow[rr] \arrow[dr] \arrow[dr, phantom, shift left=2ex, "\text{\tiny cocart.}"] & & Y_1 \arrow[dr] & & \\
 & \tilde{Z}'_0 \arrow[rr] \arrow[dr, "u_{Z_0}"'] \arrow[from=uu, crossing over] & & \tilde{Z}'_1 \arrow[dr, "u_{Z_1}"'] & \\
 & & Y'_0 \arrow[rr] \arrow[from=uuul, crossing over, "f'_0" near end] \arrow[ur, phantom, "\text{\tiny cocart.}"] & & Y'_1 \ ,
\end{tikzcd} % ?? uncertain: arrow to which the label f'_0 belongs and exact positions of the two "cocart." labels -- f'_0 printed beside the long diagonal X'_0 -> Y'_0 near Z̃'_0 -> Z̃'_1
\tag*{$(**)$}
\]
\textit{in which the cube is a totally cocartesian cube, deduced from the cocartesian square $X$ by ``exponentiating'' it by the arrow $f_0 : X_0 \longrightarrow Y_0$} \texttt{[rather, by the arrow} $(f_0,f_1)$ \texttt{of} $\underline{\Fl}(\M)$\texttt{]}\textit{.}

The base of the cube, and the square
\msp{page 33}
composed of this one and of $(*)$ (base of the polyhedral diagram $(**)$) being cartesian, the same holds for $(*)$.

\textbf{Lemma 2.7.} \textit{Let $(I_\alpha)_{\alpha\leq\alpha_0}$ be an increasing ordinal sequence of open subsets of $I$. Then $u_{I_0,I_{\alpha_0}} : Y_{I_0,I_{\alpha_0}} \longrightarrow Y_{I_{\alpha_0}}$ is a transfinite composite, indexed by the set of ordinals $-1 \leq \alpha < \alpha_0$, of a sequence of arrows $v_\alpha$, where each $v_\alpha$ is deduced from $u_{I_\alpha,I_{\alpha+1}}$ by cobase change (and from $u_{I_0} : X_{I_0} \longrightarrow Y_{I_0}$ for $\alpha = 0$).}

To see this, we introduce the transfinite sequence of objects $Z_\alpha = Y_{I_\alpha,I_{\alpha_0}}$, defined by the diagram
% foot: Version 9 février 2019

\end{document}
